\section*{第 \ref{ch:g11:functional-groups} 章——官能团与有机化合物类别}

\begin{solution}{exo:g11:functional-groups:1}
羟基，醇；链中的羰基，酮；羧基，
羧酸；酯基，酯；氨基 \ce{-NH2}，胺。
\end{solution}

\begin{solution}{exo:g11:functional-groups:2}
甲醇；丙-1-醇；甲酸；丙醛；丙酮。
\end{solution}

\begin{solution}{exo:g11:functional-groups:3}
醇；醛；酯；酰胺；卤代烷；烯烃。
\end{solution}

\begin{solution}{exo:g11:functional-groups:4}
丙醛 \ce{CH3-CH2-CHO}，丙酮 \ce{CH3-CO-CH3}。丙醛属
醛：其 \ce{C=O} 碳位于链端，与
一个氢原子相连；在丙酮中，它位于两个碳原子之间。
\end{solution}

\begin{solution}{exo:g11:functional-groups:5}
伯（连有 \ce{-OH} 的碳原子仅与一个碳原子相邻）；
仲（两个）；叔（三个）。
\end{solution}

\begin{solution}{exo:g11:functional-groups:6}
\setchemfig{atom sep=1.5em}
戊-3-醇 \chemfig{-[:30]-[:-30](-[:-90]OH)-[:30]-[:-30]};
2-甲基丁酸 \chemfig{-[:30]-[:-30](-[:-90])-[:30](=[:90]O)-[:-30]OH};
己-2-酮 \chemfig{-[:30](=[:90]O)-[:-30]-[:30]-[:-30]-[:30]};
乙酸丙酯 \chemfig{-[:30](=[:90]O)-[:-30]O-[:30]-[:-30]-[:30]}.
\end{solution}

\begin{solution}{exo:g11:functional-groups:7}
\ce{C3H6O}。是，它们互为同分异构体，但
官能团不同：一种为醛，一种为酮（二者都含羰基，但位置不同）。
\end{solution}

\begin{solution}{exo:g11:functional-groups:8}
\ce{HCOO-CH2-CH3}，甲酸乙酯；\ce{CH3-COO-CH2-CH2-CH3}，
乙酸丙酯。
\end{solution}

\begin{solution}{exo:g11:functional-groups:9}
乙酸甲酯（乙酸、甲醇）；丙酸乙酯（丙酸、
乙醇）；甲酸甲酯（甲酸、甲醇）。
\end{solution}

\begin{solution}{exo:g11:functional-groups:10}
\setchemfig{atom sep=1.5em}
\begin{center}
\small
\begin{tabular}{cccc}
\chemfig{-[:30]-[:-30]-[:30]-[:-30]OH} &
\chemfig{-[:30](-[:90]OH)-[:-30]-[:30]} &
\chemfig{-[:30](-[:90])-[:-30]-[:30]OH} &
\chemfig{-[:0](-[:90]OH)(-[:-90])-[:0]} \\
丁-1-醇 & 丁-2-醇 & 2-甲基丙-1-醇 & 2-甲基丙-2-醇 \\
伯 & 仲 & 伯 & 叔
\end{tabular}
\end{center}
\end{solution}

\begin{solution}{exo:g11:functional-groups:11}
为酰胺（乙酰胺）。\emph{-烯}。羧酸（\ce{C=O} 和
\ce{-OH} 连在同一碳原子上）、酯（\ce{C=O} 和 \ce{-O-}）及酰胺
（\ce{C=O} 和 \ce{N}）：三类。
\end{solution}

\begin{solution}{exo:g11:functional-groups:12}
阿司匹林：含一个羧基（羧酸）和一个酯基。
对乙酰氨基酚：环上有一个酚羟基 \ce{-OH}，并含酰胺基
\ce{NH-CO}。
\end{solution}

\begin{solution}{exo:g11:functional-groups:13}
\ce{C2H6O}。只有乙醇因含 \ce{O-H}，其分子间能
形成氢键（甲醚的氧原子上没有氢）。乙醇沸点更高，
为 \qty{78}{\celsius}；甲醚为
\qty{-25}{\celsius}。
\end{solution}

\begin{solution}{exo:g11:functional-groups:14}
含一个羟基和一个羧基。该链含三个碳原子，
羧基碳为1号碳，\ce{-OH} 位于2号碳上：
2-羟基丙酸。甘氨酸：2-氨基乙酸（常写作
氨基乙酸）。
\end{solution}

\begin{solution}{exo:g11:functional-groups:15}
乙醇 \ce{CH3-CH2-OH}、\ce{C2H6O}，醇；乙醛 \ce{CH3-CHO}，
\ce{C2H4O}，醛；乙酸 \ce{CH3-COOH}、\ce{C2H4O2}，
羧酸。由乙醇转化为乙醛时失去两个氢原子；
由乙醛转化为乙酸时得到一个氧原子。
\end{solution}

\begin{solution}{pb:g11:functional-groups:1}
\textbf{1.} A：酯基；B：羟基；C：羧基；D：链端的
羰基；E：酯基。

\textbf{2.} A 酯，B 醇，C 羧酸，D 醛，E 酯。

\textbf{3.} A（香蕉）与 E（菠萝）：果香气味。

\textbf{4.} 伯醇：连有 \ce{-OH} 的碳原子只与一个碳原子相连。

\textbf{5.} \ce{C6H12O}；例如己-2-酮，
\ce{CH3-CO-CH2-CH2-CH2-CH3}。

\textbf{6.} 含 \ce{C-OH} 的碳链有四个碳原子，
从 \ce{-OH} 一端编号，3 号碳上连有一个甲基：3-甲基丁-1-醇。

\textbf{7.} C：乙酸；D：己醛。

\textbf{8.} 丁酸乙酯，与丁酸和乙醇相关。

\textbf{9.} \emph{乙酸根}：酸的部分为 \ce{CH3-CO}，来自
乙酸；\emph{3-甲基丁基}：醇的部分，来自 3-甲基丁-1-醇。

\textbf{10.} C（乙酸）与 B（3-甲基丁-1-醇）。

\textbf{11.} A \ce{C7H14O2}；B \ce{C5H12O}；C \ce{C2H4O2}。

\textbf{12.} $\ce{C2H4O2} + \ce{C5H12O}$ 合计为 \ce{C7H16O3}，
多出的是一个 \ce{H2O}，即水：去掉它便得 \ce{C7H14O2}。

\textbf{13.} \ce{C2H4O2 + C5H12O -> C7H14O2 + H2O}：C 7 = 7，H 16 = 16，
O 3 = 3。

\textbf{14.} $M(\text{B}) = 5 \times 12.0 + 12 \times 1.0 + 16.0 =
\qty{88.0}{g/mol}$; $M(\text{C}) = \qty{60.0}{g/mol}$.

\textbf{15.} $60.0 + 88.0 - 18.0 = \qty{130.0}{g/mol}$.

\textbf{16.} $7 \times 12.0 + 14 \times 1.0 + 2 \times 16.0 =
\qty{130.0}{g/mol}$：两者相同。

\textbf{17.} E，\ce{C6H12O2}：$6 \times 12.0 + 12 \times 1.0 + 2 \times
16.0 = \qty{116.0}{g/mol}$.

\textbf{18.} \ce{C7H14O2}：与 A 分子式相同，因此是同分异构体，且
摩尔质量相同。仅凭摩尔质量无法区分同分异构体。

\textbf{19.} A 与庚酸原子种类和数目相同，但
官能团不同，一个是酯，一个是羧酸：一种有香蕉味，另一种
带有酸败味。官能团及其所在位置起决定作用。

\textbf{20.} \qty{130.0}{g/mol}.
\end{solution}
